\ifdefined\WPassFiveZeroNineEightLoaded\else
\gdef\WPassFiveZeroNineEightLoaded{1}
\subsection*{Passes 5098--5101: exact root-coset incidence behind the formal derivative}
The four-sheet derivative count of Passes 5090--5097 admits a stronger incidence interpretation.  For $q=2$ and $q=3$, the hypergraph whose vertices are the $q^4$ apartments through a fixed chamber and whose hyperedges are the $4q^3$ active chamber-star charts is explicitly isomorphic to the hypergraph of the maximal unipotent group $U$ of type $C_2$ and the right cosets of its four positive-root subgroups.  Thus the correspondence matches the $q$-element sheets themselves, not only their number.

At $q=3$ the derivative point graph has full automorphism group
\[
 \operatorname{Aut}(\Gamma_{\partial})\cong U_{81}\rtimes V_4,
 \qquad |\operatorname{Aut}(\Gamma_{\partial})|=324.
\]
The left-regular $U_{81}$ is normal, the identity stabilizer has order census $1^1 2^3$, and its $V_4$ fixes each of the four positive-root parallel classes setwise.  Since $324=|PGSp(4,3)|/160$, this is exactly the full projective chamber stabilizer rather than an accidental graph-symmetry enlargement.

More generally, a split finite Chevalley group with $N$ positive roots has maximal-unipotent size $q^N$ and $Nq^{N-1}$ positive-root subgroup cosets, the formal derivative of $q^N$.  In type $C_2$, the only nonzero brackets among a standard set of four positive-root nilpotents are
\[
 [X_0,X_1]=X_2,\qquad [X_0,X_2]=2X_3.
\]
Hence in odd characteristic the six unordered direction pairs generate subgroup sizes $q^2,q^2,q^2,q^2,q^3,q^4$.  Characteristic two kills the second bracket coefficient and produces the observed bad-characteristic collapse.

\paragraph{Boundary.}
The root-coset derivative law explains the chamber equality geometry and its symmetry.  It remains distinct from the still-open global expansion statement that every nontrivial Levi cohomology class activates at least $4q^3$ charts.
\fi
\ifdefined\WPassFiveOneZeroTwoLoaded\else
\gdef\WPassFiveOneZeroTwoLoaded{1}
\subsection*{Passes 5102--5109: $q=4$ expansion closure, the $U_{81}$ controller, and rank-two root calculus}
The active-chart distance attack now closes completely at $q=4$.  Two chamber stars at gallery distance $d$ share exactly $q^{4-d}$ apartments, so any XOR of $1\le m\le q$ distinct chamber stars obeys
\[
 \mathrm{wt}\ge q^3m(q+1-m)\ge q^4,
\]
with equality realized by panel subsets.  At $q=5$ this forces any counterexample to have chamber-generator leader at least six; the minimum-local-cut propagator independently forces any nonstar candidate to contain a heavier local cut.

For $q=4$, Pass5090 already proved that a word whose active $K_5$ charts are all minimum $1|4$ cuts is a chamber star.  Conditioning a heavy $2|3$ chart in each of the four fixed-apartment chart roles gives exact MILP optimum $384$.  The local $PGL(2,4)$ edge stabilizer is transitive on the six heavy cut patterns.  Hence the complete $[13600,256,256]_2$ minimum shell is exactly the $425$ chamber stars, and every nonzero word activates at least
\[
 A(y)\ge256=4q^3.
\]

The factor-$17$ bridge also becomes integral and global.  On each quadratic lane the theta companion
\[
 C=\begin{pmatrix}0&16\\1&2\end{pmatrix}
\]
and twisted block $T_6=\bigl(\begin{smallmatrix}2&8\\2&0\end{smallmatrix}\bigr)$ satisfy $CP=PT_6$ for $P=\bigl(\begin{smallmatrix}2&-2\\0&1\end{smallmatrix}\bigr)$ of determinant $2$.  Every nonsingular integral solution has even determinant, so index two is minimal.  Across the fifteen Levi line-kernel lanes this gives an explicit index-$2^{15}$ map into the exact $30$-dimensional global theta carrier.

The type-$C_2$ maximal unipotent controller now literally contains the two BT865 torsors:
\[
 U_{81}=H_{27}\,\mathbb F_3^3,\qquad
 H_{27}\cap\mathbb F_3^3=U'\cong C_3^2,
\]
with $Z(U)=Z(H_{27})\cong C_3$.  The canonical $V_4$ complement normalizes both subgroups.  On protected memory,
\[
 H_1(\mathbb F_3)\!\downarrow_U\cong\mathbb F_3[U],
\]
one regular copy; restricting to either index-three subgroup recovers the previous $3\,\mathrm{Reg}(H_{27})$ and $3\,\mathrm{Reg}(\mathbb F_3^3)$ coordinate systems.

The derivative/curvature calculus extends across rank two.  In good characteristic the pair-generated positive-root counts are encoded by
\[
 A_2:2z^2+z^3,\qquad C_2:4z^2+z^3+z^4,\qquad
 G_2:10z^2+3z^3+z^5+z^6,
\]
while the first-root derivatives are $3q^2$, $4q^3$, and $6q^5$.  The $G_2$ bad-prime reductions recorded here are only Lie-bracket shadows, not full bad-characteristic group laws.

Three auxiliary probes sharpen the same structure.  The $81\times108$ $q=3$ root-coset incidence matrix has Smith form $1^{68},3,0^{12}$, so its cokernel is $\mathbb Z^{12}\oplus\mathbb Z/3$.  The regular $U_{81}$ memory has augmentation-layer Hilbert series
\[
 (1+t+t^2)^2(1+t^2+t^4)(1+t^3+t^6),
\]
which records the $C_2$ root heights $1,1,2,3$ and restricts to the Heisenberg and flat BT865 memory filtrations.  Finally, the canonical diagonal $V_4$ splits the rational $12$-dimensional incidence defect into character sectors $4+4+2+2$.  The same defect is phase-pure for the canonical center $Z(U)=Z(H_{27})\cong C_3$: a central generator has characteristic polynomial $(x^2+x+1)^6$, so over $\mathbb C$ the defect is $6\omega\oplus6\omega^2$ with no center-trivial vectors.  Two $V_4$ involutions invert the center and exchange those conjugate six-spaces, while the third fixes the center.

\paragraph{Boundary.}
The all-$q$ and $q=5$ minimum-distance theorem remains open.  The $q=4$ minimum shell and active-chart theorem are closed.  The controller, memory, Smith-defect, and center-phase statements are exact finite group/module results, not fabricated-hardware performance claims.
\fi
%
